% #############################################################################
% This is Appendix D
% !TEX root = ../main.tex
% #############################################################################
\chapter{External Model Comparisons}
\label{chap:appendix-external-model-comparison}

This appendix evaluates four external pretrained or instruction-following
systems on the same EP/BP identification and rewriting tasks to provide
context for the adapted T5Gemma 2 models evaluated in \Cref{chap:evaluation}.
The systems are
AMALIA-9B-0626-SFT and AMALIA-9B-0626-DPO
\citep{simplicio-etal-2026-amalia}, Qwen3-4B
\citep{yang2025qwen3technicalreport}, and FLAN-T5-XL
\citep{chung2022scalinginstructionfinetunedlanguagemodels}.

\section{Evaluation Protocol}

The evaluated checkpoints are \path{amalia-llm/AMALIA-9B-0626-SFT},
\texttt{amalia-\allowbreak{}llm/\allowbreak{}AMALIA-\allowbreak{}9B-\allowbreak{}0626-\allowbreak{}DPO}, \path{Qwen/Qwen3-4B}, and
\path{google/flan-t5-xl}.

Rewriting prompts specify whether the sentence should be rewritten from BP to
EP or from EP to BP. The generated sentences are evaluated with the same BLEU
and TER scoring procedure used for the adapted models, with the two directions
reported separately.

Identification is performed differently across the external systems. AMALIA SFT
and AMALIA DPO generate an answer freely and the predicted variety label is
then parsed from the output, whereas Qwen3-4B and FLAN-T5-XL score the possible
natural-language labels directly and select the higher-scoring option. The
comparison therefore reflects each system together with its identification
inference procedure.

The systems differ in architecture, scale, training history, prompting, and
inference procedure, so the comparison provides external context rather than a
controlled model comparison.

\section{Results}

\Cref{tab:appendix-external-model-comparison} reports the external-model
results on both evaluation resources.

\begin{table}[htbp]
\centering
\thesistablefont
\renewcommand{\arraystretch}{1.08}
\setlength{\tabcolsep}{4.5pt}
\caption{External-model results on FRMT and the Golden Collection.}
\label{tab:appendix-external-model-comparison}
\begin{tabular}{@{}llccc@{}}
\toprule
\textbf{Model} & \textbf{Dataset} & \shortstack{\textbf{Macro-F1}\\\textbf{(\%)}} &
\textbf{BP-to-EP} & \textbf{EP-to-BP} \\
& & & \shortstack{\textbf{BLEU (0--100)}\\\textbf{/ TER}} & \shortstack{\textbf{BLEU (0--100)}\\\textbf{/ TER}} \\
\midrule
AMALIA SFT & FRMT & 45.37 & \textbf{44.65}/\textbf{0.4289} & 40.22/0.4682 \\
& Golden Collection & 43.99 & 68.44/0.2347 & 74.66/0.1920 \\
\addlinespace[0.2em]
AMALIA DPO & FRMT & \textbf{59.72} & 43.82/0.4494 & 35.16/0.6063 \\
& Golden Collection & \textbf{56.55} & 58.08/0.3239 & 57.18/0.3773 \\
\addlinespace[0.2em]
Qwen3-4B & FRMT & 52.56 & 39.44/0.4709 & \textbf{40.73}/\textbf{0.4654} \\
& Golden Collection & 49.96 & \textbf{78.33}/\textbf{0.1540} & \textbf{78.26}/\textbf{0.1590} \\
\addlinespace[0.2em]
FLAN-T5-XL & FRMT & 47.87 & 32.22/0.5263 & 32.17/0.5373 \\
& Golden Collection & 40.99 & 67.06/0.2209 & 66.31/0.2292 \\
\bottomrule
\end{tabular}
\end{table}

AMALIA DPO gives the highest external identification macro-F1 on both evaluation
resources, but remains below the best adapted configurations reported in
\Cref{tab:r48-classification-breakdown}.

On the Golden Collection, Qwen3-4B performs best among the external rewriting
systems in both directions, but the adapted configurations perform better for
both BP-to-EP and EP-to-BP in
\Cref{tab:r48-translation-breakdown}.

\looseness=-1
On FRMT, AMALIA SFT slightly outperforms the best adapted BP-to-EP result, with
44.65 BLEU compared with 43.97, while Qwen3-4B remains below the best adapted
EP-to-BP result.

Overall, the adapted models perform better on identification and Golden
Collection rewriting, while the external systems are competitive on FRMT only
in the BP-to-EP direction.
